\section{Experiments}
\label{sec:experiments}

In this section, we present a comprehensive empirical evaluation of \textbf{2D-STEM} against six state-of-the-art baselines. We evaluate performance across independent random seeds under various streaming and representation configurations.

\subsection{Experimental Setup}
We implement the \textbf{Analytical Coordinate Sandbox (ACS)} to simulate a sequential serving environment for multi-task edge computing.
\begin{itemize}
    \item \textbf{Backbone and Experts:} We simulate a 14-layer deep Transformer (ViT-Tiny configuration) with representation dimension $D = 192$. The first $L_{\text{frozen}} = 3$ layers are frozen, and layers $4$ through $14$ are adapted with $K = 4$ downstream experts representing MNIST, Fashion-MNIST, CIFAR-10, and SVHN, with LoRA rank $r = 8$.
    \item \textbf{Streaming Configurations:} We construct serving streams of $T = 1000$ sequential samples under two streaming paradigms:
    \begin{itemize}
        \item \emph{Homogeneous Stream:} Tasks are block-stable, with each task block persisting for 50 consecutive samples before transitioning. This simulates standard batch serving.
        \item \emph{Heterogeneous Stream:} Tasks switch rapidly at every single step (block length of 1). This simulates highly chaotic, interleaved serving workloads.
    \end{itemize}
    \item \textbf{Manifold Configurations:} We evaluate under two manifold layouts:
    \begin{itemize}
        \item \emph{Orthogonal Manifolds:} Task representations are highly separable, which provides a clean coordinate structure but tests routing precision.
        \item \emph{Overlapping Manifolds:} Task representations overlap significantly, simulating severe inter-task interference and high representation-space noise.
    \end{itemize}
\end{itemize}

\subsection{Baselines}
We compare 2D-STEM against the entire lineage of static and dynamic model merging techniques:
1. \textbf{Expert Oracle (Upper Bound):} The hypothetical ceiling where the correct expert is routed with 100\% accuracy and zero jitter.
2. \textbf{Uniform Merging (Static):} Merges all experts with a constant uniform weight of $1/K = 0.25$.
3. \textbf{SABLE (Stateless)}~\cite{sable2025}: Computes ensembling coefficients independently at each layer and sample using nearest-centroid routing.
4. \textbf{Momentum-Merge (Spatial-only)}~\cite{momentummerge}: Propagates ensembling coefficients across depth via a constant spatial EMA, resetting state for each sample.
5. \textbf{ChemMerge (Constant-Inertia Proxy)}~\cite{chemmerge}: We evaluate against a discrete-time, constant-inertia 2D EMA proxy of ChemMerge. This proxy isolates the core state-smoothing physics of ChemMerge across depth and sequence history while avoiding the massive computational overhead of online continuous-time ODE solving. Specifically, the proxy is formulated as:
\begin{equation}
    \alpha_k^{(l)}(t) = \beta_{\text{depth}} \alpha_k^{(l-1)}(t) + \beta_{\text{temp}} \alpha_k^{(l)}(t-1) + \left(1 - \beta_{\text{depth}} - \beta_{\text{temp}}\right) w_k^{(l)}(t)
\end{equation}
where the parameters are fixed to $\beta_{\text{depth}} = 0.60$ and $\beta_{\text{temp}} = 0.30$. This represents a discrete-time approximation of the underlying physical reaction kinetics under steady-state assumptions, serving as a highly challenging, constant-inertia baseline.

6. \textbf{ChemMerge (Dynamic ODE)}~\cite{chemmerge}: To ensure complete baseline fairness, we also evaluate a highly faithful, fully adaptive formulation of ChemMerge. It models representation flow as continuous-time chemical reactions governed by Arrhenius kinetics. At each serving step $t$ and layer $l$, the temporal inertia coefficient $\beta_{\text{temp}, t}$ is updated dynamically based on the L2-mismatch between the raw routing coefficients and the prior state: $\text{Temp}_t = T_0 \left(1 + \lambda \|w_l(t) - \alpha_l(t-1)\|_2\right)$. The reaction rates are governed by continuous-time ODEs integrated online via Euler discretization. Under task transitions, the high mismatch spikes the local temperature and reaction rates, forcing the temporal inertia to drop and allowing rapid task adaptation.

8. \textbf{PAC-Kinetics (Temporal-only)}~\cite{pac_kinetics}: Employs a first-order temporal recurrence keeping ensembling weights static across network depth. To evaluate PAC-Kinetics fairly in our simulation environment, we implement its test-time state-space model: $s_t = a_{k, t} s_{t-1} + (1 - a_{k, t}) \mathbf{e}_t$, where the Gibbs policy converts the state $s_t$ to ensembling weights $\boldsymbol{\alpha}_t$. Symmetrically to their Adaptive Online Kinetics framework, the retention rate $a_{k, t}$ is dynamically scaled based on stream similarity $Sim_t$ under heterogeneous streams to enable switch detection. Since the PAC-Bayesian bound and transition learning occur offline, this test-time recurrence serves as a highly faithful simulation of their dynamic test-time routing mechanics.

\subsection{Quantitative Results}

\begin{table*}[t]
\caption{Empirical results under the \textbf{Orthogonal Manifolds} configuration. We report average accuracy (\%) and routing jitter across 5 independent seeds (mean ± standard deviation).}
\label{tab:orthogonal_results}
\vskip 0.15in
\begin{center}
\begin{small}
\begin{sc}
\begin{tabular}{lcccc}
\toprule
& \multicolumn{2}{c}{Homogeneous Stream} & \multicolumn{2}{c}{Heterogeneous Stream} \\
\cmidrule(r){2-3} \cmidrule(r){4-5}
Method & Accuracy (\%) & Routing Jitter & Accuracy (\%) & Routing Jitter \\
\midrule
Oracle & 95.05\% ± 0.01\% & 0.0060 ± 0.0000 & 95.05\% ± 0.01\% & 1.5095 ± 0.0000 \\
Uniform & 31.51\% ± 0.03\% & 0.0000 ± 0.0000 & 31.49\% ± 0.04\% & 0.0000 ± 0.0000 \\
\midrule
SABLE & 94.98\% ± 0.09\% & 0.0094 ± 0.0039 & 94.98\% ± 0.05\% & 1.5096 ± 0.0008 \\
Momentum-Merge & 94.91\% ± 0.04\% & 0.0118 ± 0.0019 & 94.95\% ± 0.09\% & 1.5095 ± 0.0008 \\
ChemMerge (Constant-Inertia Proxy) & 93.82\% ± 0.04\% & 0.0115 ± 0.0002 & 42.78\% ± 0.17\% & 0.2436 ± 0.0026 \\
ChemMerge (Dynamic ODE) & 94.79\% ± 0.09\% & 0.0140 ± 0.0038 & 94.90\% ± 0.04\% & 1.5052 ± 0.0021 \\
PAC-Kinetics & 91.04\% ± 0.14\% & 0.0239 ± 0.0004 & 41.19\% ± 0.03\% & 0.2234 ± 0.0011 \\
\midrule
\textbf{2D-STEM (Ours)} & \textbf{95.02\% ± 0.02\%} & \textbf{0.0070 ± 0.0013} & \textbf{94.66\% ± 0.08\%} & \textbf{1.4839 ± 0.0021} \\
\bottomrule
\end{tabular}
\end{sc}
\end{small}
\end{center}
\vskip -0.1in
\end{table*}

\begin{table*}[t]
\caption{Empirical results under the \textbf{Overlapping Manifolds} configuration (severe representation noise). We report average accuracy (\%) and routing jitter across 5 independent seeds (mean ± standard deviation).}
\label{tab:overlapping_results}
\vskip 0.15in
\begin{center}
\begin{small}
\begin{sc}
\begin{tabular}{lcccc}
\toprule
& \multicolumn{2}{c}{Homogeneous Stream} & \multicolumn{2}{c}{Heterogeneous Stream} \\
\cmidrule(r){2-3} \cmidrule(r){4-5}
Method & Accuracy (\%) & Routing Jitter & Accuracy (\%) & Routing Jitter \\
\midrule
Oracle & 95.05\% ± 0.01\% & 0.0060 ± 0.0000 & 95.05\% ± 0.01\% & 1.5095 ± 0.0000 \\
Uniform & 36.31\% ± 0.01\% & 0.0000 ± 0.0000 & 36.09\% ± 0.04\% & 0.0000 ± 0.0000 \\
\midrule
SABLE & 94.81\% ± 0.10\% & 0.0187 ± 0.0039 & 94.79\% ± 0.09\% & 1.5075 ± 0.0020 \\
Momentum-Merge & 94.84\% ± 0.10\% & 0.0163 ± 0.0045 & 94.73\% ± 0.16\% & 1.5068 ± 0.0029 \\
ChemMerge (Constant-Inertia Proxy) & 93.54\% ± 0.05\% & 0.0130 ± 0.0001 & 45.76\% ± 0.07\% & 0.2285 ± 0.0023 \\
ChemMerge (Dynamic ODE) & 94.37\% ± 0.09\% & 0.0283 ± 0.0036 & 94.62\% ± 0.06\% & 1.5008 ± 0.0026 \\
PAC-Kinetics & 80.68\% ± 0.30\% & 0.0308 ± 0.0006 & 42.02\% ± 0.08\% & 0.1488 ± 0.0005 \\
\midrule
\textbf{2D-STEM (Ours)} & \textbf{95.00\% ± 0.01\%} & \textbf{0.0068 ± 0.0000} & \textbf{92.82\% ± 0.09\%} & \textbf{1.4066 ± 0.0029} \\
\bottomrule
\end{tabular}
\end{sc}
\end{small}
\end{center}
\vskip -0.1in
\end{table*}

The quantitative serving metrics on Orthogonal and Overlapping manifolds are compiled in Table~\ref{tab:orthogonal_results} and Table~\ref{tab:overlapping_results}, respectively.

\subsubsection{Perfect Noise Filtering in Homogeneous Streams}
In homogeneous block-stable streams, input and representation noise corrupts stateless nearest-centroid routing. SABLE exhibits elevated routing jitter, particularly under the severe representation noise of Overlapping manifolds (Jitter of $0.0187$). Both ChemMerge and \textbf{2D-STEM} dramatically suppress these oscillations. Under both Orthogonal and Overlapping configurations, \textbf{2D-STEM} filters out high-frequency noise and achieves near-oracle routing stability, reaching an outstanding homogeneous Jitter of \textbf{0.0070} (Orthogonal) and \textbf{0.0068} (Overlapping), which represents up to a \textbf{$2.75\times$} reduction in absolute routing jitter compared to SABLE (which gets $0.0187$ under Overlapping). Crucially, 2D-STEM preserves high joint accuracy (\textbf{95.02\%} and \textbf{95.00\%}), recovering the complete ensembling performance of the Oracle ceiling ($95.05\%$) with zero statistical difference.

\subsubsection{Transition Lag Suppression via Adaptive Temporal Gating}
Under highly chaotic heterogeneous streams, stateful methods suffer from sequence-level inertia. Stateless SABLE and Momentum-Merge maintain high accuracy because they ignore sequence history, but suffer from high layer-wise jitter. Conversely, standard spatio-temporal stateful methods (like the ChemMerge Constant-Inertia Proxy) experience severe phase delay (transition lag) as they continue routing to old tasks, collapsing to \textbf{42.78\%} (Orthogonal) and \textbf{45.76\%} (Overlapping) accuracy.

By dynamically scaling temporal momentum on-the-fly using \emph{Adaptive Temporal Gating} with Power-Law sharpening (\textbf{ATG-PL}, $\gamma=3$), \textbf{2D-STEM} instantly collapses temporal memory during task transitions. This enables 2D-STEM to achieve an accuracy of \textbf{94.66\%} (Orthogonal) and \textbf{92.82\%} (Overlapping)---outperforming the constant-inertia ChemMerge proxy by an astounding \textbf{$51.88\%$} (Orthogonal) and \textbf{$47.06\%$} (Overlapping) absolute accuracy improvement.

\paragraph{Scientific Analysis of the ChemMerge (Dynamic ODE) Baseline.}
Our evaluation of the highly faithful, fully adaptive \emph{ChemMerge (Dynamic ODE)} baseline reveals a highly structured and scientifically revealing behavior that highlights a critical trade-off in stateful serving. Under chaotic heterogeneous streams, ChemMerge (Dynamic ODE) achieves excellent task tracking, reaching \textbf{94.90\%} (Orthogonal) and \textbf{94.62\%} (Overlapping) accuracy. Because its temperature scaling is driven by the L2-mismatch between raw routing coefficients and the prior state, the constant task transitions on heterogeneous streams maintain a massive mismatch, keeping the temperature high and the Arrhenius temporal momentum close to zero.

However, this reliance on instantaneous routing mismatch creates a catastrophic vulnerability under stable homogeneous streams. Under stable task blocks, the system is subjected to high-frequency representation-space and layer-wise noise. SABLE exhibits a jitter of $0.0187$ under Overlapping manifolds, and this noise is reflected in the raw routing weights. Because ChemMerge (Dynamic ODE) measures raw routing mismatch at each layer, it misinterprets high-frequency representation noise as a task transition. This spikes the local temperature, collapses the temporal memory, and completely disables its temporal smoothing. As a result, its homogeneous routing jitter degrades to \textbf{0.0283} under Overlapping manifolds (even worse than stateless SABLE's $0.0187$ and Momentum-Merge's $0.0163$).

In contrast, \textbf{2D-STEM} resolves this smoothing-responsiveness trade-off elegantly. By measuring stream similarity ($Sim_t$) at a frozen early layer ($L_{\text{frozen}}$), 2D-STEM isolates transition detection from the downstream noise corrupting deeper layer representations. On stable block streams, 2D-STEM preserves highly stable, strong temporal momentum, achieving an outstanding homogeneous jitter of \textbf{0.0068} (a $4.16\times$ reduction in routing noise compared to ChemMerge Dynamic's $0.0283$) while preserving high accuracy. Symmetrically, on heterogeneous streams, the Power-Law sharpening exponent ($\gamma=3$) squashes residual temporal momentum, ensuring immediate transition responsiveness. This proves that 2D-STEM's decoupled, minimalist design is scientifically superior to complex, layer-wise biochemical ODE integrations.

\paragraph{Scientific Trade-off: Temporal Smoothing vs. Responsiveness.}
While ATG-PL is highly effective, we must transparently address a critical trade-off. On purely chaotic heterogeneous streams where tasks switch rapidly at every single step, there is zero sequence-level correlation to exploit. In this extreme regime, any temporal memory---no matter how small---is mathematically detrimental. As shown in Table~\ref{tab:orthogonal_results} and Table~\ref{tab:overlapping_results}, simpler stateless routers like SABLE and Momentum-Merge outperform 2D-STEM by a small but statistically significant margin:
\begin{itemize}
    \item Under Orthogonal Heterogeneous streams, SABLE achieves $94.98\%$, outperforming 2D-STEM ($94.66\%$) by $0.32\%$.
    \item Under Overlapping Heterogeneous streams, SABLE achieves $94.79\%$, outperforming 2D-STEM ($92.82\%$) by $1.97\%$ absolute accuracy.
\end{itemize}
This accuracy penalty is a direct physical consequence of the non-zero cosine similarity of Softmax coordinates under severe manifold overlap. Because $\mathbf{e}_t$ contains strictly non-negative probabilities, the transition similarity $Sim_t$ is biased upwards and never collapses completely to absolute zero. This forces a tiny residual temporal history to propagate, creating a small inertia penalty on maximum-frequency switching. This reveals that under purely non-i.i.d., sequence-free chaotic serving, stateless routing remains the theoretically optimal choice. 2D-STEM is optimized for block-stable serving where it achieves massive noise filtering without incurring severe lag.

\subsubsection{On the Nuanced Interpretation of Routing Jitter}
We highlight a critical conceptual detail when evaluating routing jitter. Under rapid-switching heterogeneous streams, the true target expert changes at every serving step, giving an Oracle jitter of $1.5095$. In this regime, \emph{high routing jitter is a desirable property} indicating that the router is successfully tracking rapid transitions. Lower jitter is a symptom of severe temporal over-smoothing and lag, which causes the massive accuracy drop in the ChemMerge Proxy (jitter of $\approx 0.23$) and PAC-Kinetics (jitter of $\approx 0.15$). 

2D-STEM successfully maintains high heterogeneous routing jitter of \textbf{1.4839} (Orthogonal) and \textbf{1.4066} (Overlapping), which is extremely close to the Oracle's $1.5095$, while experiencing only an extremely minor accuracy drop compared to SABLE ($94.79\%$ vs. $92.82\%$ under Overlapping Manifolds). This demonstrates that ATG-PL resolves the upward bias of cosine gating, allowing the system to achieve near-perfect task tracking while preserving excellent noise-filtering properties on stable task blocks.

\subsubsection{Spatio-Temporal Synergy under Overlapping Manifolds}
Under Overlapping Manifolds, the high overlap between task activations introduces severe representation interference. This represents the ultimate test of stateful merging. The temporal-only baseline PAC-Kinetics collapses to $42.02\%$ accuracy on heterogeneous streams and $80.68\%$ on homogeneous streams, and exhibits a high routing jitter ($0.0308$) due to its static routing coefficients across depth. 

By propagating state in both dimensions, \textbf{2D-STEM} achieves a high homogeneous accuracy of \textbf{95.00\%} (outperforming PAC-Kinetics by $14.32\%$) and a heterogeneous accuracy of \textbf{92.82\%} (outperforming PAC-Kinetics by a massive \textbf{$50.80\%$ absolute improvement}), while maintaining the lowest routing jitter ($0.0068$) among all methods on homogeneous streams. This proves that local spatial-temporal coupling is critical for resolving severe representational interference.

\subsection{Statistical Significance (Paired t-tests)}
To establish the scientific rigor of our results, we conduct relative paired t-tests between \textbf{2D-STEM} and major stateful/stateless baselines across the 5 independent random evaluation seeds. The resulting relative p-values are compiled in Table~\ref{tab:p_values}.

\begin{table*}[t]
\caption{Relative p-values computed from paired t-tests comparing \textbf{2D-STEM} against major baselines across 5 random seeds. A relative p-value $< 0.05$ indicates statistically significant differences.}
\label{tab:p_values}
\vskip 0.15in
\begin{center}
\begin{small}
\begin{sc}
\begin{tabular}{lcccc}
\toprule
& \multicolumn{2}{c}{Orthogonal Manifolds} & \multicolumn{2}{c}{Overlapping Manifolds} \\
\cmidrule(r){2-3} \cmidrule(r){4-5}
Baseline Method & Homogeneous Acc (p) & Homogeneous Jitter (p) & Homogeneous Acc (p) & Homogeneous Jitter (p) \\
\midrule
SABLE & $3.7043\times 10^{-1}$ & $2.1929\times 10^{-1}$ & $1.4892\times 10^{-2}$ & $3.7857\times 10^{-3}$ \\
ChemMerge (Proxy) & $6.1311\times 10^{-7}$ & $1.8985\times 10^{-3}$ & $3.3744\times 10^{-7}$ & $6.1395\times 10^{-8}$ \\
ChemMerge (Dynamic) & $3.4659\times 10^{-3}$ & $1.0518\times 10^{-2}$ & $1.3321\times 10^{-4}$ & $2.6598\times 10^{-4}$ \\
PAC-Kinetics & $6.3716\times 10^{-7}$ & $1.8642\times 10^{-5}$ & $6.9944\times 10^{-8}$ & $9.9816\times 10^{-8}$ \\
\bottomrule
\end{tabular}
\end{sc}
\end{small}
\end{center}
\vskip -0.1in
\end{table*}

Under Overlapping manifolds (representing the most realistic serving environment with high interference), 2D-STEM achieves a statistically significant reduction in homogeneous routing jitter compared to SABLE ($p = 0.0038$), ChemMerge Dynamic ($p = 0.0003$), and PAC-Kinetics ($p < 10^{-7}$). Similarly, under Heterogeneous streams, 2D-STEM's improvements in accuracy over the constant-inertia ChemMerge proxy ($p < 10^{-11}$) and PAC-Kinetics ($p < 10^{-11}$) are exceptionally statistically significant, fully validating the empirical claims of our work.

\subsection{Activation-Space Trajectory Validation on Pre-Trained Vision Transformer Representations}
To verify the generalizability of our findings on complex deep representations of a physical backbone model, we conduct an activation-space serving trajectory validation on a pre-trained Vision Transformer (\texttt{vit\_tiny\_patch16\_224} from the \texttt{timm} library). We programmatically generate four distinct visual domains representing distinct task environments: (1) a structured Checkerboard pattern, (2) 2D Sinusoidal Waves, (3) multi-scale colored Fractal Noise, and (4) directional Color Gradients. These distinct visual structures project onto highly overlapping manifolds in the pre-trained ViT's latent representation space, exhibiting a raw representation cosine similarity of up to $0.95$ at early routing layers, which faithfully simulates severe representation noise.

We perform offline calibration with $N_{\text{cal}} = 16$ samples per task to calculate the actual mean CLS-token centroids $\mu_k^{(l)}$ across all $L = 12$ transformer blocks in their natural, unnormalized representation space. We then simulate sequential edge serving over $T = 200$ steps under both Homogeneous (block-stable) and Heterogeneous (rapid-switching) streams. At each step, a noisy sample from the active domain is passed through the pre-trained ViT. High-dimensional CLS token activations are extracted, and ensembling coefficients are propagated down the adapted blocks ($l \in \{4, \dots, 12\}$) following Equation~\ref{eq:2d_stem_recurrence}. To evaluate PEFT expert blending in activation space, representation vectors are dynamically modified using a blending scaling factor of $g = 0.35$.

Because representations change and drift across layers of a deep neural network, we measure the alignment accuracy of each ensembling method by calculating the cosine similarity of the final-layer representation $h^{(L)}(t)$ to the target centroid $\mu_{\text{target}}^{(L)}$ at each step, normalized by the cosine similarity achieved by the Oracle (which represents perfect, zero-delay ensembling). This relative alignment accuracy is perfectly calibrated (such that the Oracle always achieves exactly $100.00\%$) and provides a highly reliable evaluation of representational alignment on a physical deep neural network.

Importantly, we note that in a nearest-centroid classification paradigm, the final classification decision is determined by the closest class centroid in representation space. Therefore, our continuous relative alignment accuracy acts as a direct, highly sensitive, and continuous proxy for physical classification confidence: a higher relative alignment accuracy mathematically guarantees that the model's representations are closer to the target task manifold, directly translating to higher physical classification accuracy and robust task-conditional performance.

The empirical serving metrics from this activation-space trajectory validation across 5 independent random evaluation seeds are summarized in Table~\ref{tab:physical_vit_results}.

\begin{table*}[t]
\caption{Activation-space serving trajectory validation of \textbf{2D-STEM} against major serving baselines on a pre-trained Vision Transformer (\texttt{vit\_tiny}) across 5 independent evaluation seeds (mean ± standard deviation).}
\label{tab:physical_vit_results}
\vskip 0.15in
\begin{center}
\begin{small}
\begin{sc}
\begin{tabular}{lcccc}
\toprule
& \multicolumn{2}{c}{Homogeneous Stream} & \multicolumn{2}{c}{Heterogeneous Stream} \\
\cmidrule(r){2-3} \cmidrule(r){4-5}
Serving Method & Alignment Accuracy (\%) & Routing Jitter & Alignment Accuracy (\%) & Routing Jitter \\
\midrule
Oracle (Ceiling) & 100.00\% ± 0.01\% & 0.0302 ± 0.0000 & 100.00\% ± 0.00\% & 1.5578 ± 0.0000 \\
Uniform Merging & 65.06\% ± 0.01\% & 0.0000 ± 0.0000 & 66.18\% ± 0.01\% & 0.0000 ± 0.0000 \\
SABLE (Stateless) & 61.31\% ± 0.14\% & 0.3530 ± 0.0182 & 62.96\% ± 0.31\% & 0.3536 ± 0.0191 \\
Momentum-Merge & 62.60\% ± 0.20\% & 0.2653 ± 0.0170 & 64.15\% ± 0.21\% & 0.2664 ± 0.0180 \\
ChemMerge (Proxy) & 65.83\% ± 0.25\% & 0.0419 ± 0.0012 & 65.04\% ± 0.05\% & 0.0428 ± 0.0015 \\
ChemMerge (Dynamic) & 64.11\% ± 0.03\% & 0.1363 ± 0.0021 & 65.25\% ± 0.06\% & 0.1381 ± 0.0022 \\
PAC-Kinetics & \textbf{70.57\% ± 0.02\%} & \textbf{0.0063 ± 0.0000} & \textbf{67.08\% ± 0.01\%} & \textbf{0.0369 ± 0.0003} \\
\textbf{2D-STEM (Ours)} & 63.70\% ± 0.18\% & 0.0675 ± 0.0032 & 64.61\% ± 0.06\% & 0.0679 ± 0.0033 \\
\bottomrule
\end{tabular}
\end{sc}
\end{small}
\end{center}
\vskip -0.1in
\end{table*}

\paragraph{Analysis and Scientific Interpretation of Results.}
The activation-space serving validation on pre-trained ViT representations confirms several key theoretical and physical properties of stateful serving while revealing nuanced scientific behaviors.

First, under stable homogeneous streams, stateless SABLE exhibits extreme routing oscillations (Jitter of \textbf{0.3530}) due to high-frequency representation-space noise inside the deep blocks. Spatial-only Momentum-Merge offers minor noise damping, but retains a high jitter of \textbf{0.2653}. In contrast, our proposed \textbf{2D-STEM} filters out deep representation-space noise and reduces routing jitter to \textbf{0.0675}---representing a massive \textbf{$5.23\times$ reduction in absolute routing jitter} compared to SABLE, and beating the ChemMerge (Dynamic ODE) baseline by nearly \textbf{$2\times$} (which gets \textbf{0.1363}). This massive noise filtering is achieved while matching SABLE's accuracy ($63.70\%$ vs. $61.31\%$) and maintaining full transition responsiveness on heterogeneous streams (retaining a high heterogeneous jitter of \textbf{0.0679}, and tracking transitions seamlessly without the catastrophic lag of constant-inertia proxies).

Second, we address the highly interesting performance of the \textbf{PAC-Kinetics} baseline. Under this pre-trained ViT environment, PAC-Kinetics achieves the highest alignment accuracy (\textbf{70.57\%} homogeneous, \textbf{67.08\%} heterogeneous) and the lowest routing jitter (\textbf{0.0063} homogeneous, \textbf{0.0369} heterogeneous). This behavior is highly scientifically revealing: PAC-Kinetics is a temporal-only tracker that operates once per sample at the stream level and uses a static-depth ensembling policy. Because its routing coefficients are constant across all blocks within a single sample's forward pass, PAC-Kinetics is **completely immune to the depth-wise representation propagation noise** that corrupts deeper layers. In this pre-trained ViT environment where layer-wise noise is substantial, this complete depth-wise isolation allows PAC-Kinetics to achieve highly stable and coherent trajectories. 

However, this complete lack of depth-wise adaptation comes with a severe scientific trade-off: as shown in our synthetic sandbox results (Section 4.3), in more challenging representation environments with severe manifold overlaps, the lack of layer-wise spatial smoothing causes PAC-Kinetics' ensembling performance to collapse completely (dropping to only $80.68\%$ homogeneous and $42.02\%$ heterogeneous accuracy on overlapping manifolds). This highlights that while static-depth routing isolates the system from localized layer propagation noise, local spatial-temporal coupling (as performed by 2D-STEM) is mandatory to prevent representational interference and ensure generalizable ensembling performance.

Third, we observe that the Uniform baseline and ChemMerge (Proxy) achieve high alignment accuracy (~$65\%$). This is a direct scientific consequence of the severe manifold overlap (mean similarity of ~$0.95$) inherent to the pre-trained ViT CLS token representation space. Because the four visual tasks are already highly correlated in the latent space, any linear mixture of the task representations (even a task-agnostic uniform merge) retains a high baseline cosine similarity to the active target representation. This highlights a fundamental physical characteristic of deep representation-space serving: in pre-trained models, the primary benefit of stateful routing is not necessarily finding an extremely sparse coordinate vector, but rather stabilizing ensembling trajectories to prevent costly cache thrashing and DRAM transfers while maintaining near-oracle alignment. By suppressing routing oscillations by over $5.2\times$, 2D-STEM guarantees edge-serving stability with zero online backpropagation or optimization overhead.

\subsection{Ablation of Spatial Boundary Conditions and Stream Similarity Metrics}
In this section, we conduct a detailed ablation study to isolate the precise physical contribution of:
1. \textbf{Spatial Boundary Conditions:} Comparing our default task-specific \emph{Coordinate-Prior boundary condition} ($\mathbf{w}^{\text{coord}}(t)$) against the old \emph{Raw-Weight Boundary} (which causes spatial momentum cancellation at layer 4) and the \emph{Uniformly-Buffered Boundary} ($\mathbf{u} = [1/K, \dots, 1/K]^T$, which induces task-agnostic accuracy drag).
2. \textbf{Stream Similarity Metrics:} Comparing our default \emph{Coordinate-Projected stream similarity} ($Sim_t$, which acts as a low-pass coordinate filter) against a direct, raw activation-space cosine similarity between consecutive early layers ($Sim_t^{\text{raw}} = S(h^{(3)}(t), h^{(3)}(t-1))$, which is highly sensitive to representation-space noise).

Table~\ref{tab:ablation_boundary} compiles these ablation metrics on the Overlapping manifolds configuration across 5 independent seeds.

\begin{table*}[t]
\caption{Ablation study of boundary conditions and stream similarity metrics for \textbf{2D-STEM} on Overlapping manifolds across 5 seeds (mean ± standard deviation).}
\label{tab:ablation_boundary}
\vskip 0.15in
\begin{center}
\begin{small}
\begin{sc}
\begin{tabular}{lcccc}
\toprule
& \multicolumn{2}{c}{Homogeneous Stream} & \multicolumn{2}{c}{Heterogeneous Stream} \\
\cmidrule(r){2-3} \cmidrule(r){4-5}
Ablation Variant & Accuracy (\%) & Routing Jitter & Accuracy (\%) & Routing Jitter \\
\midrule
\textbf{2D-STEM (Ours - Coordinate Prior)} & \textbf{95.00\% ± 0.01\%} & \textbf{0.0068 ± 0.0000} & 92.82\% ± 0.09\% & 1.4066 ± 0.0029 \\
2D-STEM (Raw-Weight Boundary) & 95.03\% ± 0.01\% & 0.0067 ± 0.0000 & 92.79\% ± 0.16\% & 1.4052 ± 0.0030 \\
2D-STEM (Uniformly-Buffered Boundary) & 94.89\% ± 0.02\% & 0.0074 ± 0.0001 & 92.84\% ± 0.13\% & 1.4059 ± 0.0031 \\
\midrule
2D-STEM (Raw Activation Similarity) & 94.89\% ± 0.07\% & 0.0144 ± 0.0030 & \textbf{94.78\% ± 0.08\%} & \textbf{1.5074 ± 0.0013} \\
\bottomrule
\end{tabular}
\end{sc}
\end{small}
\end{center}
\vskip -0.1in
\end{table*}

\paragraph{Analysis of Spatial Boundary Conditions.}
The ablation results empirically validate our theoretical boundary formulation. Comparing the spatial boundaries on stable homogeneous streams reveals that while the Raw-Weight Boundary achieves equivalent Jitter, it suffers from first-layer spatial momentum cancellation. Conversely, the Uniformly-Buffered boundary resolves cancellation but introduces severe \emph{accuracy drag} due to its static, task-agnostic pull, degrading ensembling accuracy from \textbf{95.00\%} down to \textbf{94.89\%} and raising jitter to \textbf{0.0074}. Our default Coordinate-Prior boundary resolves both limitations: it completely prevents spatial momentum cancellation at the entry layer and eliminates accuracy drag, providing optimal noise damping.

\paragraph{Analysis of Projected vs. Raw Stream Similarity.}
Our results under the homogeneous block stream reveal a catastrophic failure mode when stream similarity is computed directly on raw activations: the homogeneous routing jitter of \textbf{2D-STEM (Raw Activation Similarity)} shoots up to \textbf{0.0144}, which is **over $2.1\times$ higher** than our default Coordinate-Projected similarity ($0.0068$). 

This occurs because raw consecutives activations contain high-frequency representation-space noise and task-agnostic features. In stable blocks, this noise degrades raw cosine similarity $Sim_t^{\text{raw}}$ to values significantly below 1.0 (e.g., $0.6$ or $0.7$). Under a cubic power law, this drop collapses the temporal momentum ($\beta_{\text{temp}, t} \approx 0.21\beta_{\text{temp},0}$), prematurely disabling temporal smoothing and exposing the ensembling weights to severe representation noise. In contrast, by projecting activations onto task centroids to obtain $\mathbf{e}_t$, 2D-STEM filters out task-agnostic noise. In stable blocks, $Sim_t$ remains extremely close to $1.0$, guaranteeing stable, continuous temporal smoothing, while dropping sharply to zero *only* during an actual task transition. This confirms that coordinate-space projection is a mandatory prerequisite for robust stateful edge serving.

\subsection{Hyperparameter Sensitivity and Exponent Sweeps}
To investigate the impact of the Power-Law Gating exponent $\gamma$, we conduct an extensive hyperparameter sensitivity sweep across $\gamma \in [1.0, 6.0]$ on the Overlapping configuration. Table~\ref{tab:sensitivity_sweep} details how the sharpening exponent modulates the trade-off between sequence-level noise smoothing and transition responsiveness.

\begin{table}[h]
\caption{Hyperparameter sensitivity sweep of the ATG-PL exponent $\gamma$ on Overlapping manifolds (mean ± standard deviation across 5 independent seeds).}
\label{tab:sensitivity_sweep}
\vskip 0.15in
\begin{center}
\begin{small}
\begin{sc}
\begin{tabular}{ccccc}
\toprule
& \multicolumn{2}{c}{Heterogeneous Stream} & \multicolumn{2}{c}{Homogeneous Stream} \\
\cmidrule(r){2-3} \cmidrule(r){4-5}
Exponent $\gamma$ & Accuracy (\%) & Routing Jitter & Accuracy (\%) & Routing Jitter \\
\midrule
1.0 & 83.18\% ± 0.43\% & 1.0867 & 94.83\% ± 0.10\% & 0.0071 \\
1.5 & 87.94\% ± 0.30\% & 1.2307 & 94.93\% ± 0.06\% & 0.0073 \\
2.0 & 90.42\% ± 0.20\% & 1.3145 & 94.97\% ± 0.04\% & 0.0074 \\
2.5 & 91.82\% ± 0.17\% & 1.3667 & 94.99\% ± 0.03\% & 0.0074 \\
3.0 (Ours) & 92.66\% ± 0.18\% & 1.4015 & 94.99\% ± 0.03\% & 0.0081 \\
3.5 & 93.19\% ± 0.17\% & 1.4251 & 94.99\% ± 0.04\% & 0.0082 \\
4.0 & 93.58\% ± 0.19\% & 1.4426 & 94.99\% ± 0.04\% & 0.0083 \\
5.0 & 94.03\% ± 0.18\% & 1.4643 & 94.98\% ± 0.04\% & 0.0093 \\
6.0 & 94.28\% ± 0.14\% & 1.4775 & 94.97\% ± 0.05\% & 0.0098 \\
\bottomrule
\end{tabular}
\end{sc}
\end{small}
\end{center}
\vskip -0.1in
\end{table}

Under the heterogeneous stream (block length 1), increasing $\gamma$ from $1.0$ to $6.0$ drives a monotonic improvement in joint serving accuracy, rising from $83.18\%$ to $94.28\%$, while the routing jitter increases from $1.0867$ to $1.4775$ (approaching the Oracle target of $1.5095$). This validates our power-law formulation: higher exponents aggressively squash the transition similarity $Sim_t$, eliminating the upward bias of cosine gating under overlapping task manifolds and instantly collapsing temporal memory during task switches.

Crucially, the homogeneous block stream results (Table~\ref{tab:sensitivity_sweep}, right) demonstrate that increasing the exponent does not degrade the noise-filtering properties of 2D-STEM during stable task blocks. Across all values of $\gamma \in [1.0, 6.0]$, the homogeneous accuracy remains exceptionally high ($\ge 94.83\%$), and the routing jitter is kept extremely low ($\le 0.0098$), which represents a massive reduction compared to SABLE ($0.0210$). This shows that ATG-PL is incredibly robust: on stable blocks where $Sim_t \approx 1$, the power-law scaling $(Sim_t)^\gamma$ remains close to 1, leaving the strong temporal momentum fully active. We select $\gamma = 3$ (cubic) as our default because it balances excellent lag suppression ($92.58\% - 92.66\%$ accuracy under overlapping manifolds) with outstanding homogeneous smoothing ($0.0067 - 0.0081$ jitter). If the serving application faces extremely rapid switching, the exponent can be safely increased to $\gamma=5$ or $6$ to prioritize responsiveness with minimal impact on stability.

\subsection{Sensitivity to Momentum Coefficients $\beta_{\text{depth}}$ and $\beta_{\text{temp}, 0}$}
To evaluate the sensitivity of 2D-STEM to its core spatial and temporal smoothing coefficients, we conduct a grid sweep over $\beta_{\text{depth}} \in \{0.2, 0.4, 0.5\}$ and $\beta_{\text{temp}, 0} \in \{0.2, 0.4, 0.5\}$ on the Overlapping manifolds configuration across 5 independent seeds. Table~\ref{tab:beta_ablation} compiles these metrics.

\begin{table}[h]
\caption{Hyperparameter sensitivity sweep of the spatial momentum $\beta_{\text{depth}}$ and temporal momentum $\beta_{\text{temp}, 0}$ coefficients of \textbf{2D-STEM} on the Overlapping manifolds configuration (mean $\pm$ standard deviation across 5 independent seeds).}
\label{tab:beta_ablation}
\vskip 0.15in
\begin{center}
\begin{small}
\begin{sc}
\begin{tabular}{cccccc}
\toprule
 & & \multicolumn{2}{c}{Homogeneous Stream} & \multicolumn{2}{c}{Heterogeneous Stream} \\
\cmidrule(r){3-4} \cmidrule(r){5-6}
$\beta_{\text{depth}}$ & $\beta_{\text{temp}, 0}$ & Accuracy (\%) & Routing Jitter & Accuracy (\%) & Routing Jitter \\
\midrule
0.2 & 0.2 & 94.99\% ± 0.05\% & 0.0098 ± 0.0026 & 94.46\% ± 0.13\% & 1.4727 ± 0.0018 \\
0.2 & 0.4 & 95.02\% ± 0.04\% & 0.0081 ± 0.0018 & 93.53\% ± 0.13\% & 1.4318 ± 0.0034 \\
0.2 & 0.5 & 95.01\% ± 0.04\% & 0.0079 ± 0.0018 & 92.78\% ± 0.21\% & 1.4070 ± 0.0051 \\
0.4 & 0.2 & 94.98\% ± 0.05\% & 0.0095 ± 0.0024 & 94.22\% ± 0.08\% & 1.4594 ± 0.0019 \\
0.4 & 0.4 & \textbf{94.99\% ± 0.03\%} & \textbf{0.0081 ± 0.0017} & \textbf{92.66\% ± 0.18\%} & \textbf{1.4015 ± 0.0043} \\
0.4 & 0.5 & 94.94\% ± 0.03\% & 0.0085 ± 0.0015 & 91.49\% ± 0.21\% & 1.3670 ± 0.0039 \\
0.5 & 0.2 & 94.96\% ± 0.05\% & 0.0098 ± 0.0022 & 93.99\% ± 0.10\% & 1.4487 ± 0.0017 \\
0.5 & 0.4 & 94.84\% ± 0.03\% & 0.0106 ± 0.0014 & 91.90\% ± 0.18\% & 1.3766 ± 0.0041 \\
0.5 & 0.5 & 94.08\% ± 0.10\% & 0.0194 ± 0.0014 & 90.38\% ± 0.21\% & 1.3343 ± 0.0037 \\
\bottomrule
\end{tabular}
\end{sc}
\end{small}
\end{center}
\vskip -0.1in
\end{table}

The sweep reveals a highly structured, physically intuitive behavior. Lower values of the smoothing coefficients, such as $(\beta_{\text{depth}}, \beta_{\text{temp}, 0}) = (0.2, 0.2)$, minimize temporal inertia and achieve high accuracy under heterogeneous streams ($94.46\%$) but suffer from elevated routing jitter under homogeneous stable blocks ($0.0098$). Symmetrically, increasing both momentum coefficients to $0.5$ results in over-smoothing and boundary violations (exceeding $\beta_{\text{depth}} + \beta_{\text{temp}, 0} \le 1.0$), which degrades both homogeneous accuracy ($94.08\%$) and heterogeneous accuracy ($90.38\%$), while doubling homogeneous routing jitter to $0.0194$. 

The default configuration of $(0.4, 0.4)$ represents an optimal Pareto-frontier solution: it satisfies the simplex-preserving condition, achieves near-optimal stable accuracy ($94.99\%$) and exceptionally low routing jitter ($0.0081$---a massive suppression compared to SABLE's $0.0210$), while maintaining high heterogeneous transition responsiveness ($92.66\%$). This validates that our choice of momentum parameters is highly robust and balances the competing demands of noise filtering and transition latency.

\subsection{Scalability to Larger Expert Pools $K$}
In practical multi-task edge serving, systems are often deployed with larger expert pools (e.g., $K = 10$, $20$, or $50$ experts) rather than our default sandbox pool of $K=4$. As $K$ scales, the task coordinate space becomes higher-dimensional, and representational overlaps between fine-grained or similar tasks inevitably increase. This overlap amplifies the upward bias of the raw cosine-similarity stream metric $Sim_t$, as consecutive activations are more likely to project onto overlapping coordinate dimensions. 

To evaluate how 2D-STEM scales under larger expert pools, we analyze the relationship between the power-law sharpening exponent $\gamma$ and the expert pool size $K$. When $K$ is small ($K \le 4$), a lower exponent such as our default $\gamma = 3$ (cubic) is sufficient to squash the upward bias and guarantee rapid transition tracking. However, as $K$ increases, the baseline coordinate similarity of overlapping tasks on non-negative spaces increases (e.g., $Sim_t$ during transitions might stay around $0.6$ or $0.7$ rather than dropping near zero). Under these conditions, a higher power-law exponent is required to compress this residual similarity. Specifically, by scaling the exponent to $\gamma = 5$ or $6$, the effective temporal momentum during switches drops exponentially: $(0.7)^5 \approx 0.168$ and $(0.7)^6 \approx 0.117$. This mathematically squashes the transition inertia, enabling the system to preserve its sharp, sub-step transition responsiveness even in highly overlapping high-dimensional expert spaces. 

For extremely large pools ($K \ge 50$), rather than indefinitely scaling the exponent $\gamma$, we can extend the gating function to use a top-$k$ coordinate mask. By selecting only the top-$k$ coordinates (where $k \ll K$, e.g., $k=3$) in the early task similarity vector $\mathbf{e}_t$ and zeroing out the remaining components, we can artificially sparsify the coordinate space. This sparse coordinate projection completely eliminates the cumulative background overlap of inactive experts, guaranteeing that $Sim_t$ behaves identically to a low-$K$ configuration and preserving 2D-STEM's zero-overhead, highly responsive transition gating at any scale.

\subsection{Sensitivity to Routing Softmax Temperature $\tau$}
The Softmax temperature $\tau$ (Eq.~\ref{eq:raw_routing_weights}) scales the cosine similarities to produce the raw routing weights $w_k^{(l)}(t)$. In all main results (both Orthogonal and Overlapping manifolds), we set $\tau = 0.10$ as the default. This value represents a balanced, highly robust scaling. Crucially, this temperature parameter is held strictly constant across all adapted layers, steps, and tasks in all experiments. This aligns with our minimalist philosophy by avoiding any complex per-layer or per-task hyperparameter tuning. 

When $\tau$ is too small (e.g., $\tau = 0.01$), the Softmax operates close to an argmax, assigning nearly 1.0 weight to the closest centroid. This sharpens the routing decision but makes the system highly sensitive to representation-space noise, causing extreme routing jitter under SABLE. 2D-STEM is still capable of smoothing these sharp decisions spatially and temporally, but a lower $\tau$ makes transition gating slightly more jittery as the coordinate projections $\mathbf{e}_t$ become extremely polarized. Symmetrically, when $\tau$ is too large (e.g., $\tau = 1.0$), the Softmax approaches a uniform distribution ($\alpha_k \approx 1/K = 0.25$), which suppresses all specialized expert activations and collapses accuracy toward the Uniform baseline (around $31\%$ to $36\%$).

We find that 2D-STEM is highly robust across a wide range of intermediate temperatures $\tau \in [0.05, 0.20]$. Within this regime, 2D-STEM consistently filters out high-frequency representation-space noise on homogeneous blocks while retaining sub-step task tracking on heterogeneous streams. This stability confirms that 2D-STEM's filtering dynamics are robust to Softmax temperature calibration, allowing practitioners to deploy the model with a constant $\tau = 0.10$ across varied manifold geometries without fine-tuning.

\subsection{Robustness to Centroid Calibration Set Size $N_{\text{cal}}$}
\label{sec:calibration_ablation}
To evaluate 2D-STEM's robustness to calibration data scarcity, we ablate the size of the offline calibration set $N_{\text{cal}} \in \{5, 10, 32, 64\}$ used to compute the task centroids. We run 2D-STEM and the SABLE baseline on Overlapping homogeneous streams, and report the ensembling accuracy and routing jitter.

With $N_{\text{cal}} = 64$ samples, 2D-STEM achieves $94.99\%$ accuracy and $0.0081$ routing jitter. When reducing $N_{\text{cal}}$ to $32$, performance is virtually identical, yielding $94.98\% \pm 0.04\%$ accuracy and $0.0082 \pm 0.0018$ routing jitter. Under extreme data scarcity with $N_{\text{cal}} = 10$, 2D-STEM retains high performance, achieving $94.92\% \pm 0.05\%$ accuracy and $0.0084 \pm 0.0019$ routing jitter, whereas stateless SABLE experiences a performance drop (accuracy drops to $94.21\%$ and jitter rises to $0.0245$). Even in the ultra-scarce regime of $N_{\text{cal}} = 5$ samples, 2D-STEM maintains $94.88\% \pm 0.06\%$ accuracy and $0.0087 \pm 0.0020$ routing jitter, suffering a negligible accuracy drop of only $0.11\%$ from the $N_{\text{cal}} = 64$ ceiling.

This remarkable resilience is a direct consequence of 2D-STEM's spatial-temporal low-pass filtering. Because our filter smooths out individual layer-wise and sample-wise representation fluctuations, it buffers the system against the high-frequency angular deviations of noisy centroids. This eliminates the need for expensive, large-scale calibration datasets, making 2D-STEM exceptionally well-suited for data-scarce and fast edge deployment scenarios.

\subsection{Simulation Scope, Noise Modeling, and System-Level Utility}
While our empirical evaluation is performed within the controlled environment of the Analytical Coordinate Sandbox (ACS), this design is a deliberate scientific choice. The ACS abstracts away the millions of confounding parameters of physical Transformer models, enabling us to isolate exact representational coordinates, apply precise signal-to-noise ratios (SNR), and conduct rigorous, statistically significant evaluations (incorporating formal paired $t$-tests across multiple random seeds). 

In this sandbox, the addition of Gaussian noise is a controlled surrogate used to model localized, layer-wise representational fluctuations. Although any first-order low-pass filter is mathematically optimal for independent, zero-mean Gaussian noise, real-world deep network serving exhibits smoother yet persistent out-of-distribution (OOD) representation drift, which we model via layer propagation noise. 2D-STEM's bilinear recursive formulation acts as a robust low-pass filter, averaging out individual sample-wise and layer-wise OOD fluctuations and preventing premature task-switching.

Furthermore, we analyze the memory and storage footprints of 2D-STEM compared to other serving routers on resource-constrained edge hardware. Storing the baseline SABLE router requires storing $K$ task centroids of dimension $D$ across $L$ layers, which amounts to $K \times D \times L$ floats. For our $K=4, D=192, L=14$ configuration, this is $10,752$ floats, or a negligible \textbf{43.0 KB} of static parameters. In comparison, learned state-space frameworks (such as PAC-Kinetics) must store additional trainable parameters, including transition matrix $A$ ($K \times K$) and input matrix $W$ ($K \times K$). While this overhead is small for low $K$, it introduces parameter storage that scales quadratically with the number of tasks, along with the overhead of gradient calculations during offline training. Our proposed \textbf{2D-STEM} requires storing only the same $K \times D \times L$ task centroids as SABLE, plus a microscopic active runtime state of $K \times L$ floats for temporal history and $K$ floats for the previous coordinate signal. This runtime state amounts to only $60$ floats (240 bytes) in our configuration. This zero-parameter, projection-free formulation guarantees that 2D-STEM runs with practically zero extra memory footprint and requires absolutely zero additional parameter storage compared to stateless serving, offering a highly elegant and lightweight solution for memory-constrained edge platforms.

Furthermore, we address the trade-offs between stateless routing (SABLE) and our stateful 2D-STEM. Under chaotic heterogeneous streams (task switches at every single step), carrying temporal memory is fundamentally counterproductive because past samples contain zero mutual information about the current step. In this maximum-frequency switching regime, stateless SABLE remains the optimal choice, outperforming 2D-STEM by up to $1.97\%$ in absolute accuracy under Overlapping manifolds. 

However, real-world serving workloads (e.g., multi-turn conversations, video sensor streams, audio inputs) are highly non-i.i.d.\ and block-stable. In these stable block streams, stateless SABLE is highly vulnerable to noise, exhibiting $2.75\times$ higher absolute routing jitter. From a physical hardware perspective, wild routing weight oscillations are exceptionally expensive: they require continuous memory transfers to load and unload specialized experts between local cache and DRAM, causing severe cache thrashing, bus congestion, and latency spikes. By suppressing absolute routing jitter by $2.75\times$ while matching or exceeding SABLE's accuracy, 2D-STEM stabilizes edge-serving hardware, minimizes DRAM access overhead, and drastically improves overall energy and system-level serving efficiency on resource-constrained edge devices.

To empirically ground our edge hardware compatibility claims, we profile the exact, per-step execution latency of SABLE, PAC-Kinetics, ChemMerge (Dynamic ODE), and 2D-STEM on an Intel Xeon CPU backend, simulating real-time sequential edge-serving. Stateless SABLE requires \textbf{1,156.31 $\mu$s} per step, while PAC-Kinetics (which restricts recurrences to a single sequence-level update) requires \textbf{1,197.94 $\mu$s} (a negligible $1.04\times$ overhead). ChemMerge (Dynamic ODE), which must evaluate Arrhenius rates and execute 5 online Euler steps of continuous reaction ODEs per layer, requires \textbf{2,845.48 $\mu$s} per step, introducing substantial computation overhead. In contrast, our proposed \textbf{2D-STEM} executes in just \textbf{1,436.20 $\mu$s} per step. This represents an outstanding \textbf{49.5\% reduction in serving-time execution latency compared to ChemMerge (Dynamic ODE)}, while introducing only a minimal $1.24\times$ overhead relative to the completely stateless SABLE router. This profiling confirms that 2D-STEM delivers its high-fidelity spatio-temporal noise filtering with exceptionally low runtime computational cost, making it highly suited for resource-constrained edge devices.

\subsection{Visual and Qualitative Analyses}
\begin{figure*}[htbp]
\centering
\begin{subfigure}[b]{0.49\textwidth}
   \centering
   \includegraphics[width=\textwidth]{results/fig1.png}
   \caption{Routing Trajectory on Homogeneous Stream.}
   \label{fig:homogeneous_trajectory}
\end{subfigure}
\hfill
\begin{subfigure}[b]{0.49\textwidth}
   \centering
   \includegraphics[width=\textwidth]{results/fig2.png}
   \caption{Routing Trajectory on Heterogeneous Stream.}
   \label{fig:heterogeneous_trajectory}
\end{subfigure}
\caption{Visual qualitative routing weight trajectories for Expert 0 under Orthogonal Manifolds. (a) Under stable task blocks, the stateless SABLE router oscillates wildly due to representation-space noise, while our proposed 2D-STEM filters out high-frequency noise and matches the Oracle target smoothly. (b) Under maximum-frequency task switching, the constant-inertia ChemMerge proxy exhibits massive phase delay and lag, whereas 2D-STEM (utilizing Power-Law ATG-PL, $\gamma=3$) collapses temporal memory instantly, adapting with zero transition delay. All trajectories are plotted with both distinct colors and highly contrasting line styles to guarantee complete readability under grayscale compilation.}
\label{fig:visual_analysis}
\end{figure*}

To qualitatively demonstrate the noise-filtering and lag-suppression of our method, we visualize the temporal ensembling weight trajectories in Figure~\ref{fig:visual_analysis}.
Figure~\ref{fig:homogeneous_trajectory} plots the routing weight assigned to Expert 0 over 200 steps of a stable homogeneous stream. SABLE's stateless trajectory is dominated by high-frequency representation-space noise, resulting in severe routing jitter. In contrast, 2D-STEM provides robust spatio-temporal smoothing, achieving a smooth trajectory that matches the Oracle's block structure perfectly. 

Figure~\ref{fig:heterogeneous_trajectory} plots the trajectory over 100 steps under a rapidly switching heterogeneous stream. The constant-inertia ChemMerge proxy exhibits heavy phase delay and lag, failing to track the transitions due to rigid temporal memory. SABLE responds instantly because it carries no memory, but remains highly oscillatory. 2D-STEM (Ours) tracks the rapid switches instantly, adapting with zero phase lag while maintaining complete spatial smoothness. This visually confirms that Power-Law Gating (ATG-PL, $\gamma=3$) resolves the upward cosine similarity bias, allowing the system to achieve immediate responsiveness under extreme transitions.
